\documentclass[10pt,a4paper]{amsart}
\usepackage[margin=19mm]{geometry}
\usepackage{amsmath,amssymb,mathtools}
\usepackage[scr=rsfs,cal=euler]{mathalfa}
\usepackage{xfrac}
\usepackage[unicode,hidelinks,bookmarks=false]{hyperref}
\newcommand{\ie}{i.e.\ }
\newcommand{\cf}{cf.\ }
\newcommand{\resp}{respectively\ }
\newcommand{\Subsection}{\subsection}
\providecommand{\nobreakdash}{\nobreak\hbox{-}}
\setlength{\emergencystretch}{2em}
\multlinegap0pt
\usepackage{amssymb}
\usepackage{xcolor}
\newtheorem{lemma}{Lemma}
\newtheorem{corollary}{Corollary}
\newtheorem{proposition}{Proposition}
\DeclareMathOperator{\Div}{div}
\newcommand{\Kmat}{\mathsf K}
\newcommand{\Real}{\mathbb R}
\newcommand{\Rn}{\Real^n}
\newcommand{\Ufield}{\mathsf U}
\newcommand{\breg}{\mathfrak b}
\DeclareMathOperator{\curl}{curl}
\newcommand{\eframe}{\mathsf e}   % Cartesian frame along X
\newcommand{\gmetric}{\mathbf g}
\newcommand{\preg}{\mathfrak p}
\newcommand{\sreg}{\mathfrak s}
\DeclareMathOperator{\tr}{tr}
\title{Ideal MHD below the classical well-posedness threshold}
\author{Matteo Giardi}
\date{Source: arXiv:2609.01093v1 (CC BY 4.0)}
\begin{document}
\maketitle

\noindent\emph{Source and licence.} Ideal MHD below the classical well-posedness threshold. Version arXiv:2609.01093v1 (CC BY 4.0), distributed by arXiv under the Creative Commons Attribution 4.0 International licence, http://creativecommons.org/licenses/by/4.0/, as recorded in the arXiv licence metadata for this exact version and shown on the abstract page https://arxiv.org/abs/2609.01093v1. Full licence notice: \texttt{licenses/arxiv-2609.01093-CC-BY-4.0.txt}.\par\smallskip

\noindent\textit{Editorial scope.} Counted unit: the article's proposition prop:constraint (source0.tex lines 3377--3386). The article's complete original proof of it is delivered with it (source0.tex lines 3388--3596, one \texttt{proof} environment of 209 lines). No mathematics is written, repaired or re-derived: the statement and the proof are the article's own lines, in the article's own notation, and no retained line is substituted or re-flowed.  Retained uncounted as the source-local context the proof needs: lines 1095--1097; lines 1145--1150; lines 1152--1185; lines 1259--1280; lines 1286--1337; lines 1606--1618; lines 2052--2062.  Nothing was omitted from any retained span, so the package carries no omission record.  No retained line contains a citation, so the wrapper carries no bibliography and no bibliography entry.  No retained line contains a figure environment, an image-inclusion command or drawing code, so no image is delivered and the package declares no asset dependency.  Editorial formatting only, in addition: an AMS \texttt{amsart} preamble that restates the article's own declarations and macros used by the retained lines, namely the environments \texttt{lemma}, \texttt{corollary}, \texttt{proposition} on separate counters, together with the standard packages those lines need.  The retained environments print in this document as Lemma 1, Corollary 1, Proposition 1 in order of appearance, whereas the article prints the same items with the numbering of its own full document.  The complete native source of the article as distributed in the arXiv e-print for this exact version is bundled unchanged as \texttt{sources/209079/source0.tex}.
\begin{equation}\label{eq:metric-tensor-divergence}
 (\Div_{\gmetric}T)^j:=E_iT^{ji}.
\end{equation}

\begin{equation}\label{eq:invdiv}
    \begin{cases}
        \curl_{\gmetric}F=0,\\
     \Div_{\gmetric}F=q.
\end{cases}
\end{equation}

\begin{lemma}[Canonical solvability of the metric Hodge system]
\label{lem:hodge-solvability}
Let $X:\Real_y^n\to\Real_x^n$ be an orientation-preserving global
bi-Lipschitz map with $d\mu=J dy$.  Fix $1<p<n$ and put
\[
 p^*:=\frac{np}{n-p}.
\]
For every $q\in L^p(\mu)$ there is a unique vector field $F=F^i\eframe_i$ along $X$, whose Cartesian components belong to
$L^{p^*}(\mu;\Rn)$, satisfying \eqref{eq:invdiv} in the conjugated
distributional sense: equivalently, $u:=F\circ X^{-1}$ satisfies
\[
 \partial_{x^i}u^j-\partial_{x^j}u^i=0,
 \qquad
 \partial_{x^i}u^i=q\circ X^{-1}
 \quad\text{in }\mathcal D'(\Real_x^n).
\]
The solution is
\begin{equation}\label{eq:pressure-gradient}
 \mathcal P_{\gmetric}q
 :=\mathcal C_X\nabla\Delta_x^{-1}\mathcal C_X^{-1}q, \qquad \mathcal{C}_Xf=f\circ X,
\end{equation}
and it satisfies
\begin{equation}\label{eq:canonical-hls}
 \|\mathcal P_{\gmetric}q\|_{L^{p^*}(\mu)}
 \le C_{n,p}\|q\|_{L^p(\mu)}.
\end{equation}
In particular, for
\begin{equation}\label{eq:rHLS}
    r_{\mathrm{HLS}}:=\frac{4n}{n+4},
\end{equation}
one has $1<r_{\mathrm{HLS}}<n$ for every $n\ge2$ and $r_{\mathrm{HLS}}^*=4$; hence
every $q\in L^{r_{\mathrm{HLS}}}(\mu)$ has a unique solution
$\mathcal P_{\gmetric}q\in L^4(\mu)$ which we refer to as the canonical Hodge vector field along $X$.
\end{lemma}

If $\Ufield=dw$, the off-constraint auxiliary equation considered below, without imposing $J[\Ufield]=1$, is
the wave--Hodge system
\begin{subequations}\label{eq:auxiliary-system}
\begin{align}
&
\begin{cases}
    \Box w=-F[dw],\\
    (w,\partial_tw)|_{t=0}=(0,v_0),
\end{cases}\\
 &
 \begin{cases}
  \curl_{\gmetric[dw]}F[dw]&=0,\\
  \Div_{\gmetric[dw]}F[dw]&=\mathcal Q[dw]
 \end{cases}
 &\label{eq:force-hodge}
\end{align}
\end{subequations}
The last two equations, together with the normalisation
$F[dw]\in L^4(\mu[dw])$, are precisely the defining characterization of
$F[dw]$ furnished by Lemma~\ref{lem:hodge-solvability}; without the
integrability killing the harmonic components, they determine $F[dw]$ only up to $(\nabla H)\circ X$ with $H$
harmonic.\\

\begin{lemma}[weak pullback calculus]\label{lem:weak-pullback}
Let $X$ be an orientation-preserving global bi-Lipschitz map as above.
\begin{enumerate}
\item[(i)] $E_i$ is skew-adjoint on $L^2(\mu)$. In particular
  \[
   \int_{\Real_y^n}(E_if)\phi d\mu=-\int_{\Real_y^n}f(E_i\phi)d\mu
  \]
  whenever $f,\phi,E_if,E_i\phi\in L^2(\mu)$.
\item[(ii)] Let $F=F^i\eframe_i$ be a vector field along $X$ whose Cartesian components belong to $L^1_{\rm loc}(\Real^n_y;\Rn)$. If
  $q,c^{ij}\in L^1_{\rm loc}(\mu)$, then
  \[
   \Div_{\gmetric}F=q,
   \qquad
   (\curl_{\gmetric}F)^{ij}=c^{ij}
  \]
  in the conjugated distributional sense if and only if
  \[
   \partial_a(J\Kmat_i^aF^i)=Jq,
   \qquad
   \partial_a\!\left(J\Kmat_i^aF^j-J\Kmat_j^aF^i\right)=Jc^{ij}
   \quad\text{in }\mathcal D'(\Real^n_y).
  \]
  Equivalently, for every $\varphi\in C_c^\infty(\Real^n_y)$,
  \begin{equation}\label{eq:conservativeform}
   \int q\varphi d\mu=-\int F^i\Kmat_i^a\partial_a\varphi d\mu,\qquad
   \int c^{ij}\varphi d\mu=-\int\bigl(F^j\Kmat_i^a-F^i\Kmat_j^a\bigr)
          \partial_a\varphi d\mu.
  \end{equation}
  With the convention in \eqref{eq:metric-tensor-divergence}, the same
  calculation gives, whenever
  $H:=\Div_{\gmetric}T\in L^1_{\rm loc}(\mu)$,
  \begin{equation}\label{eq:conservative-tensor-divergence}
   JH^j
   =\partial_a(J\Kmat_i^aT^{ji})
   \quad\text{in }\mathcal D'(\Real_y^n)
  \end{equation}
  for every locally integrable tensor $T$.
\item[(iii)] If $\mathcal C_X^{-1}F\in W^{1,p}_{\rm loc}$ for some
  $1\le p\le\infty$, then 
  \begin{equation}\label{eq:biLipchainrule}
      \partial_aF^i=A_a^jE_jF^i,
  \end{equation}
  and the conjugated divergence and curl in
  (ii) agree almost everywhere with the pointwise formulas $\Kmat^a_i\partial_aF^i$ and
  $\Kmat^a_i\partial_aF^j-\Kmat^a_j\partial_aF^i$.
\item[(iv)] The conjugated derivatives commute on their common
distributional domain:
  \[
   E_iE_j=E_jE_i.
  \]
\end{enumerate}
\end{lemma}

\begin{corollary}[traces in the $Y$ scale]\label{cor:trace-Y}
For $\breg>\frac12$, $\sreg\ge0$, any $\preg$, and every
$f\in Y^{\sreg,\preg}_{\breg}$,
$$\sup_t\big\|\langle D_{\parallel}\rangle^{\preg}\Lambda^{\sreg}f(t)\big\|_{L^{2}}\ \le\
C_{\breg} \|f\|_{Y^{\sreg,\preg}_{\breg}},$$
and $t\mapsto f(t)$ is continuous into $Y^{\sreg,\preg}$. Consequently
\begin{equation}\label{eq:trace-class}
H^{\sreg,\preg}_{\breg}\ \hookrightarrow\ C(\Real;Y^{\sreg,\preg})\cap C^1(\Real;Y^{\sreg,\preg-1}),
\qquad
\sup_t\big(\|f(t)\|_{Y^{\sreg,\preg}}+\|\partial_tf(t)\|_{Y^{\sreg,\preg-1}}\big)
\le C_{\breg}\|f\|_{H^{\sreg,\preg}_{\breg}}.
\end{equation}
\end{corollary}

\begin{lemma}[uniqueness for the homogeneous wave equation]\label{lem:wave-uniqueness}
Let $0<T<\infty$ and let $u\in C^1([0,T];L^2(\Rn))$ satisfy
$\mathcal Bu\in C([0,T];L^2(\Rn))$,
\[
 \Box u=0\quad\text{in }\mathcal D'((0,T)\times\Rn),
 \qquad
 u(0)=\partial_tu(0)=0 .
\]
Then $u\equiv0$ on $[0,T]$. The same conclusion holds componentwise for
vector- and matrix-valued $u$.
\end{lemma}

\begin{proposition}[weak Jacobi's identity]\label{prop:constraint}
Let $w$ solve \eqref{eq:auxiliary-system} on $[0,T]$ with $v_0$ divergence-free and
\begin{equation}\label{eq:constraint-hyp}
\Ufield=dw\ \text{on }[0,T]\times\Rn
\ \text{for some }\Ufield\in H^{s-1,1}_{\breg},
\qquad \big\|dw\big\|_{L^\infty}\le\tfrac12,\qquad
d(\Box X)\in L^1_tL^2_{y} .
\end{equation}
Then $\det dX\equiv1$ on $[0,T]$.
\end{proposition}

\begin{proof}
Recall the notation
\[
 \sigma=s-1
 ,
 \qquad
 A=dX=\mathrm{Id}+dw,
 \qquad
 \Kmat=A^{-1},
 \qquad
 \varrho:=\log\det A.
\]
The trace estimate of Corollary~\ref{cor:trace-Y} applied to $\Ufield$, together with
$\Ufield=dw$ on $[0,T]$, gives
\[
 dw\in C_tY^{\sigma,1},
 \qquad
 \partial_tdw,\ \mathcal Bdw\in C_tH^\sigma.
\]
Consequently,
\[
 \mathcal A^\pm dw
 =\partial_tdw\pm\mathcal Bdw
 \in L^2_tH^\sigma
 \hookrightarrow L^2_tL^4_y,
\]
because $\sigma>(n-1)/2\ge n/4$.  The pointwise assumption
$\|dw\|_{L^\infty}\le 1/2$ keeps $A$ in a fixed compact subset of
the invertible matrices and therefore gives
\[
 \Kmat\in L^\infty.
\]

For smooth $X$ the following computation is immediate from Jacobi's formula.
At the present regularity, differentiating the first-order identity once more
and identifying the resulting quadratic product are not automatic, so we
justify the full second-order identity by mollification. Fix
$I'\Subset(0,T)$, mollify $dw$ in spacetime using only values inside
$(0,T)$, and set
\[
 A_\ell:=\mathrm{Id}+(dw)_\ell,
 \qquad
 \Kmat_\ell:=A_\ell^{-1},
 \qquad
 \varrho_\ell:=\log\det A_\ell.
\]
The mollifiers preserve
$\|A_\ell-\mathrm{Id}\|_{L^\infty}\le1/2$, so the matrices
$\Kmat_\ell$ are uniformly bounded. All the differential operators involved
have constant coefficients and commute with mollification. Jacobi's formula
and differentiation of $\Kmat_\ell A_\ell=\mathrm{Id}$ give
\[
 \mathcal A^\pm\varrho_\ell
 =\tr\bigl(\Kmat_\ell\mathcal A^\pm A_\ell\bigr),
 \qquad
 \mathcal A^+\Kmat_\ell
 =-\Kmat_\ell(\mathcal A^+A_\ell)\Kmat_\ell.
\]
Therefore
\begin{align*}
 \Box\varrho_\ell
 &=\mathcal A^+\tr\bigl(\Kmat_\ell\mathcal A^-A_\ell\bigr)\\
 &=\tr\bigl(\Kmat_\ell d(\Box X)_\ell\bigr)
   -\tr\bigl(
       \Kmat_\ell(\mathcal A^+A_\ell)
       \Kmat_\ell(\mathcal A^-A_\ell)
     \bigr).
\end{align*}
Here $\mathcal A^\pm A_\ell=\mathcal A^\pm(dw)_\ell$ and
$\Box A_\ell=d(\Box X)_\ell$. We use the cyclicity of the trace explicitly:
\[
 \tr\bigl(
   \Kmat_\ell\mathcal A^+(dw)_\ell
   \Kmat_\ell\mathcal A^-(dw)_\ell
 \bigr)
 =\tr\bigl(
   \Kmat_\ell\mathcal A^-(dw)_\ell
   \Kmat_\ell\mathcal A^+(dw)_\ell
 \bigr).
\]
After averaging these equal expressions, the definition of $Q_0$ gives
\begin{align*}
 \tr(\Kmat_\ell \mathcal A^+(dw)_\ell \Kmat_\ell \mathcal A^-(dw)_\ell)
 &=\frac12\tr\bigl(
      \Kmat_\ell \mathcal A^+(dw)_\ell \Kmat_\ell \mathcal A^-(dw)_\ell
      +\Kmat_\ell \mathcal A^-(dw)_\ell \Kmat_\ell \mathcal A^+(dw)_\ell
    \bigr)\\
 &=\frac12\Kmat_{\ell,i}^a\Kmat_{\ell,j}^b
   \Big[
      (\mathcal A^+(dw)_\ell)_a^j
      (\mathcal A^-(dw)_\ell)_b^i
     +(\mathcal A^-(dw)_\ell)_a^j
      (\mathcal A^+(dw)_\ell)_b^i
   \Big]\\
 &=\Kmat_{\ell,i}^a\Kmat_{\ell,j}^b
   Q_0\bigl((dw)_{\ell,a}^{ j},(dw)_{\ell,b}^{ i}\bigr)\\
 &=-\mathcal Q[(dw)_\ell].
\end{align*}
Hence
\begin{equation*}
    \Box\varrho_\ell
 =\mathcal Q[(dw)_\ell]
  +\tr\bigl(\Kmat_\ell d(\Box X)_\ell\bigr).
\end{equation*}
Put
\[
 P_\ell^\pm:=\Kmat_\ell\mathcal A^\pm(dw)_\ell,
 \qquad
 P^\pm:=\Kmat\mathcal A^\pm dw,
 \qquad
 G_\ell:=d(\Box X)_\ell,
 \qquad
 G:=d(\Box X).
\]
After passing to a subsequence, $\Kmat_\ell\to\Kmat$ almost everywhere,
and the matrices $\Kmat_\ell$ are uniformly bounded. Standard mollifier
convergence gives
$\mathcal A^\pm(dw)_\ell\to\mathcal A^\pm dw$ in
$L^2(I';L^4)$ and $G_\ell\to G$ in $L^1(I';L^2)$. The decompositions
\begin{align*}
 P_\ell^\pm-P^\pm
 &=\Kmat_\ell\bigl(\mathcal A^\pm(dw)_\ell
                    -\mathcal A^\pm dw\bigr)
   +(\Kmat_\ell-\Kmat)\mathcal A^\pm dw,\\
 \Kmat_\ell G_\ell-\Kmat G
 &=\Kmat_\ell(G_\ell-G)+(\Kmat_\ell-\Kmat)G
\end{align*}
and dominated convergence therefore give
\[
 \begin{gathered}
  \varrho_\ell\longrightarrow\varrho
  \quad\text{in }L^2_{\mathrm{loc}}(I'\times\Rn),\\
  P_\ell^\pm\longrightarrow P^\pm
  \quad\text{in }L^2(I';L^4_y),\\
  \tr(\Kmat_\ell G_\ell)\longrightarrow\tr(\Kmat G)
  \quad\text{in }L^1(I';L^2_y).
 \end{gathered}
\]
Here the first convergence follows as well from the Lipschitz continuity of
$M\mapsto\log\det(\mathrm{Id}+M)$ on the fixed compact matrix range. Since
$\mathcal Q$ is a finite sum of products of $P_\ell^+$ and $P_\ell^-$,
the $L^2_tL^4_y$ convergence gives
$\mathcal Q[(dw)_\ell]\to\mathcal Q[dw]$ in $L^1(I';L^2_y)$. Since $I'$ was
arbitrary, passing to the limit gives
\begin{equation}\label{eq:weak-jacobi}
 \Box\varrho
 =\mathcal Q[dw]+\tr\bigl(\Kmat d(\Box X)\bigr)
 \qquad\text{in }\mathcal D'((0,T)\times\Rn).
\end{equation}

Passing to the limit in the first-order Jacobi formulas also gives
\begin{equation}\label{eq:constraint-first-order-jacobi}
 \partial_t\varrho
 =\tr(\Kmat\partial_tdw),
 \qquad
 \mathcal B\varrho
 =\tr(\Kmat\mathcal Bdw).
\end{equation}
Since $dw\in C_tY^{\sigma,1}\hookrightarrow C_tL^\infty$ and the maps
$A\mapsto A^{-1}$ and $A\mapsto\log\det A$ are Lipschitz on the fixed set of
matrices under consideration, we have
\[
 \Kmat\in C([0,T];L^\infty),
 \qquad
 \varrho\in C([0,T];L^2).
\]
Together with $\partial_tdw,\mathcal Bdw\in C_tL^2$, the first-order
identities imply
\[
 \varrho\in C^1([0,T];L^2),
 \qquad
 \mathcal B\varrho\in C([0,T];L^2).
\]
In particular, they hold also at $t=0$ and $t=T$.

The trace bounds give
$\mathcal Q[dw]\in L^\infty_tL^{r_{\mathrm{HLS}}}_y$, so the
Hardy--Littlewood--Sobolev estimate and the uniform Jacobian bounds give
$F\in L^\infty_tL^4_y$. Together with
$dF=-d(\Box X)\in L^1_tL^2_y$, this gives
$F\in L^1_tW^{1,2}_{\rm loc}$. Bi-Lipschitz composition gives the same local
regularity for $\mathcal C_X^{-1}F$, so Lemma~\ref{lem:weak-pullback}
identifies $\tr(\Kmat dF)$ with $\Div_{\gmetric}F$. Using the Hodge equation
in \eqref{eq:weak-jacobi}, we conclude that
\[
 \Box\varrho
 =\mathcal Q[dw]-\Div_{\gmetric}F
 =0
 \qquad\text{in }\mathcal D'((0,T)\times\Rn).
\]
The Cauchy data for $X$ give
\[
 A(0)=\mathrm{Id},
 \qquad
 \varrho(0)=0,
\]
and \eqref{eq:constraint-first-order-jacobi} gives
\[
 \partial_t\varrho(0)
 =\tr(dv_0)
 =\Div_yv_0
 =0.
\]

Since $\varrho\in C^1([0,T];L^2)$ with $\mathcal B\varrho\in C([0,T];L^2)$,
$\Box\varrho=0$ in $\mathcal D'$ and $\varrho(0)=\partial_t\varrho(0)=0$,
Lemma~\ref{lem:wave-uniqueness} gives $\varrho\equiv0$.
Therefore $\det dX=\exp(\varrho)\equiv1$, as claimed.
\end{proof}

\end{document}
